\section{JAX Shallow-Water Solver}
\label{sec:jax-solver}

\subsection{Structured computational domain}

The JAX solver operates on cell-averaged water depth \(h\) and momenta \(hu\)
and \(hv\) on a regular Cartesian raster.  For the configured run, the domain,
synthetic DEM, and LULC roughness data are resampled to a 50.0~m grid in the
projected coordinate system of the domain.  Bounds are aligned outward to the
grid spacing.  The DEM is bilinearly resampled, with nearest-neighbour and
nearest-finite-cell fallbacks used only where the first interpolation leaves an
in-domain value undefined.  Roughness values are rasterised from the
\texttt{roughness} attribute.  Failure to assign a finite DEM or roughness
value to any active cell raises an error.

Raster rows are reversed after geospatial resampling so that both the array row
index and the exported \(y\) coordinate increase northwards.  The production
grid regression test passed with a shape of 517 by 551 cells, EPSG:32643, and
finite bed and roughness values throughout the active mask.

\subsection{Finite-volume update}

Fluxes at cell faces are evaluated with the Harten--Lax--van Leer (HLL)
approximate Riemann construction~\cite{HartenLaxVanLeer1983}.  For each
coordinate direction, the implementation estimates the bounding wave speeds
from the local normal velocity and \(\sqrt{gh}\), where the configured
gravitational acceleration is \(9.81\)~m~s\(^{-2}\).  The left or right
physical flux is selected when both waves travel in one direction; otherwise,
the two-wave HLL intermediate flux is used.

Bed variation is treated using hydrostatic reconstruction and paired
face-source corrections following the well-balanced construction of
Audusse et al.~\cite{AudusseEtAl2004}.  At each face, reconstructed depths are
formed relative to the higher adjacent bed elevation and clipped at zero.
The paired pressure corrections are then applied on the two sides of the
face.  This implementation is intended to preserve a stationary free surface
over variable bathymetry while maintaining non-negative reconstructed depth.

The update adds spatially distributed rainfall directly to water depth.
Rainfall records are required to carry \texttt{mm/hr} units and are converted
to metres per second using a factor of \(10^{-3}/3600\).  Rates are linearly
interpolated in simulation time and multiplied by the active-domain rainfall
mask.  Manning friction is applied semi-implicitly to both momentum
components after the flux and rainfall update.  Negative updated depth is
clipped to zero, and momentum is set to zero wherever depth does not exceed the
configured dry tolerance.

\subsection{Time integration and configured experiment}

The time step is selected from an unsplit two-dimensional wave-speed
restriction using a configured CFL number of 0.45 and is capped at 10.0~s.
The dry-depth tolerance is \(1.0\times10^{-4}\)~m.  The complete adaptive
integration is represented by a JIT-compiled \texttt{jax.lax.scan} with a
maximum of 20,000 steps; the implementation does not use a Python loop to
advance physical time.  Requested output times truncate individual steps so
that fields are recorded at the prescribed interval.

The configured development run uses 32-bit floating-point arithmetic, starts
from zero depth and zero velocity, and spans 3600.0~s with output every
60.0~s.  The rainfall scenario identifier is
\nolinkurl{45rp_rain}, and its default rate outside the record is
0.0~mm~hr\(^{-1}\).  The configuration notes that this one-hour window
excludes the rainfall peak and is not a complete event simulation.

\subsection{Boundary treatment}

The numerical core supports reflective and zero-gradient open ghost cells.
At a reflective external boundary, the normal momentum changes sign while
depth and tangential momentum are copied.  At an active--inactive interface
inside the rectangular raster, the active state is mirrored and its normal
momentum is reversed.  This retains the hydrostatic pressure contribution
while cancelling mass transport across the irregular domain mask.

The production configuration selects \texttt{all\_reflective}, matching the
final ANUGA policy \texttt{all\_reflective\_final}.  No open segment is
inferred from the unresolved source attributes.  This makes the two numerical
models internally comparable, but does not validate the closed-boundary
assumption against real drainage.

\subsection{Output and numerical verification}

The saved NetCDF dataset contains elevation, depth, both velocity components,
and velocity magnitude on common \((\texttt{time},\texttt{y},\texttt{x})\)
coordinates.  Out-of-domain cells are represented by missing values, and
depth and velocity are labelled in metres and metres per second,
respectively.

All JAX solver regression tests passed in the validation run.  In the
reflective flat-basin test, water volume at all requested times satisfied the
implemented comparison tolerance of relative \(2.0\times10^{-7}\) and absolute
\(1.0\times10^{-5}\) in grid-cell depth sum.  In the variable-bed
lake-at-rest test, final free-surface elevation was compared with an absolute
tolerance of \(2.0\times10^{-7}\)~m, while the maximum absolute value of each
momentum component was required to remain below \(1.0\times10^{-6}\).  A
separate irregular-mask test applied the same depth tolerance and momentum
limits at internal solid walls.  All three tests passed.

Cross-solver consistency was tested on a 10 by 10 flat basin with 10.0~m
cells, reflective boundaries, no friction, a 20.0~s duration, and a local
inflow rate of 0.005~m~s\(^{-1}\).  The JAX and ANUGA peak depths at the
inflow location were required to agree within 15 per cent relative tolerance
and 0.001~m absolute tolerance.  Arrival at a downstream test cell was defined
by a 0.001~m depth threshold and required to agree within 3.0~s.  The test
passed.  These checks detect major conservation, balance, boundary, and
cross-solver inconsistencies; they do not constitute validation against an
observed flood event.
